\begin{figure}[htbp]
\centering
\begin{tikzpicture}[x=1.05cm,y=1cm,>=Latex,font=\small]
 \draw[very thick,ink] (0,0)--(9,0);
 \draw[very thick,dashed,accent] (9,0)--(10,0);
 \foreach \x in {0,...,10}{
   \draw (\x,-.08)--(\x,.08);
   \node[below=5pt,font=\scriptsize] at (\x,0) {$\x$};
 }
 \foreach \x in {1,...,9}{
  \draw[->,ink] (\x-.85,.35)--(\x-.15,.35);
  \node[above=3pt] at (\x-.5,.35) {$\x$};
 }
 \draw[->,accent] (9.15,.35)--(9.85,.35);
 \node[above=3pt,text=accent] at (9.5,.35) {$1$};
 \draw[decorate,decoration={brace,amplitude=5pt},ink] (0,1)--(9,1)
  node[midway,above=7pt] {Nine oriented intervals: one turn};
 \node[align=center] at (5,-1.05)
  {Transition index: $g(t)=1+(t\bmod9)$\\
   Memory of the lifted calendar: $m(t)=\lfloor t/9\rfloor$};
 \node[align=center] at (5,-1.85)
  {$g(t+9)=g(t),\qquad m(t+9)=m(t)+1$};
\end{tikzpicture}
\caption{Lift of the cycle of nine transitions. The upper numbers
label oriented intervals; the lower numbers locate their boundaries.
The dashed interval shows the first transition of the next turn.
The phase index returns and the calendar quotient increases. This diagram
represents the calendar; the action on residues, quotients, orientation,
boundary, and memory is defined through the operations of the TPK.}
\label{fig:ii-segmento-generador}
\end{figure}
